\documentclass[11pt]{article}
\usepackage[utf8]{inputenc}
\usepackage[spanish, es-tabla]{babel}
\usepackage{amsmath, amssymb}
\usepackage{booktabs}
\usepackage{enumitem}
\usepackage[margin=2.5cm]{geometry}

\setlength{\parindent}{0pt}
\setlength{\parskip}{0.65em}

\title{Econometría ICS-294\\Guía de ejercicios: Clase 5}
\author{Francisco Alfaro Medina}
\date{}

\begin{document}
\maketitle

\section*{Objetivo}

Esta guía busca practicar el efecto del error de medición en regresiones lineales, distinguiendo entre errores en la variable dependiente y errores en la variable independiente.

\section*{Ejercicios}

\begin{enumerate}[label=\textbf{Ejercicio \arabic*.}, leftmargin=*]

\item \textbf{Variable verdadera y variable observada.}

Para cada caso, identifique la variable verdadera $W^*$, la variable observada $W$ y el posible error de medición $u=W-W^*$.

\begin{enumerate}[label=\alph*)]
    \item Ingreso real mensual versus ingreso declarado en una encuesta.
    \item Calorías efectivamente consumidas versus calorías compradas.
    \item Horas reales de estudio versus horas reportadas por estudiantes.
    \item Habilidad cognitiva verdadera versus puntaje en una prueba estandarizada.
\end{enumerate}

\item \textbf{Error clásico de medición.}

Explique qué significa que el error de medición sea clásico:
\[
E(u)=0,\qquad \operatorname{Cov}(u,W^*)=0.
\]

Luego indique un ejemplo donde este supuesto podría ser razonable y otro donde podría fallar.

\item \textbf{Error de medición en la variable dependiente.}

Suponga que el modelo verdadero es:
\[
Y^*=\beta_0+\beta_1X+\varepsilon,
\qquad E[\varepsilon\mid X]=0,
\]
pero solo observamos:
\[
Y=Y^*+u.
\]

\begin{enumerate}[label=\alph*)]
    \item Escriba el modelo estimado usando $Y$.
    \item ¿Cuál es el nuevo término de error?
    \item Si $\operatorname{Cov}(X,u)=0$, ¿$\hat{\beta}_1$ es sesgado?
    \item ¿Qué ocurre con la varianza del término de error?
\end{enumerate}

\item \textbf{Demostración: error en $Y$.}

Use:
\[
E(\hat{\beta}_1)=\frac{\operatorname{Cov}(X,Y)}{\operatorname{Var}(X)}
\]
y $Y=Y^*+u$ para mostrar que, si $\operatorname{Cov}(X,u)=0$, entonces:
\[
E(\hat{\beta}_1)=\beta_1.
\]

\item \textbf{Precisión estadística.}

Explique por qué el error de medición en la variable dependiente no sesga la pendiente, pero sí puede hacer más difícil encontrar efectos estadísticamente significativos.

Use la idea:
\[
\operatorname{Var}(\varepsilon+u)=\operatorname{Var}(\varepsilon)+\operatorname{Var}(u)
\]
cuando $\varepsilon$ y $u$ no están correlacionados.

\item \textbf{Error de medición en la variable independiente.}

Suponga que el modelo verdadero es:
\[
Y=\beta_0+\beta_1X^*+\epsilon,
\]
pero observamos:
\[
X=X^*+u.
\]

\begin{enumerate}[label=\alph*)]
    \item ¿Por qué este caso es más problemático que medir mal $Y$?
    \item Si se estima una regresión de $Y$ sobre $X$, ¿qué contiene el nuevo término de error?
    \item Explique intuitivamente por qué $X$ queda correlacionada con el error del modelo estimado.
\end{enumerate}

\item \textbf{Sesgo de atenuación.}

Bajo error clásico en $X$, se tiene:
\[
E(\hat{b}_1)=
\beta_1
\frac{\operatorname{Var}(X^*)}
{\operatorname{Var}(X^*)+\operatorname{Var}(u)}.
\]

\begin{enumerate}[label=\alph*)]
    \item Defina el factor de atenuación $a$.
    \item Explique por qué $0<a<1$ si $\operatorname{Var}(u)>0$.
    \item ¿Qué ocurre con $E(\hat{b}_1)$ si $\operatorname{Var}(u)$ aumenta?
    \item ¿Qué ocurre si $\operatorname{Var}(u)=0$?
\end{enumerate}

\item \textbf{Cálculo del factor de atenuación.}

Suponga que $\beta_1=4$, $\operatorname{Var}(X^*)=9$ y $\operatorname{Var}(u)=3$.

\begin{enumerate}[label=\alph*)]
    \item Calcule el factor de atenuación $a$.
    \item Calcule $E(\hat{b}_1)$.
    \item ¿Cuánto se subestima el efecto verdadero?
\end{enumerate}

\item \textbf{Signo del coeficiente.}

Suponga que el efecto verdadero es negativo: $\beta_1=-2$. Además, el factor de atenuación es $a=0{,}6$.

\begin{enumerate}[label=\alph*)]
    \item Calcule $E(\hat{b}_1)$.
    \item ¿El estimador queda más cerca o más lejos de cero?
    \item Explique por qué el sesgo de atenuación reduce la magnitud del coeficiente en valor absoluto.
\end{enumerate}

\item \textbf{Signal-to-noise ratio.}

Compare los siguientes escenarios.

\begin{center}
\begin{tabular}{ccc}
\toprule
Escenario & $\operatorname{Var}(X^*)$ & $\operatorname{Var}(u)$ \\
\midrule
A & 100 & 5 \\
B & 100 & 50 \\
C & 100 & 200 \\
\bottomrule
\end{tabular}
\end{center}

\begin{enumerate}[label=\alph*)]
    \item Calcule el factor de atenuación en cada escenario.
    \item ¿En cuál hay mayor sesgo?
    \item Interprete la importancia relativa de señal y ruido.
\end{enumerate}

\item \textbf{Salario e IQ.}

Se quiere estimar el efecto de la habilidad cognitiva verdadera sobre el salario:
\[
salario_i=\beta_0+\beta_1 habilidad_i^*+\epsilon_i.
\]
Sin embargo, solo se observa un puntaje de IQ que mide la habilidad con error.

\begin{enumerate}[label=\alph*)]
    \item ¿Cuál es la variable independiente verdadera?
    \item ¿Cuál es la variable independiente observada?
    \item ¿Qué tipo de sesgo esperaría en el coeficiente del puntaje de IQ?
    \item ¿El coeficiente estimado probablemente sobreestima o subestima la relación verdadera entre habilidad y salario?
\end{enumerate}

\item \textbf{Comparación final.}

Complete la tabla.

\begin{center}
\begin{tabular}{p{4cm}p{4cm}p{4cm}}
\toprule
Caso & ¿Sesga la pendiente? & Principal consecuencia \\
\midrule
Error clásico en $Y$ & & \\
Error clásico en $X$ & & \\
\bottomrule
\end{tabular}
\end{center}

Luego escriba un párrafo breve explicando por qué ambos errores de medición no tienen las mismas consecuencias.

\end{enumerate}

\section*{Preguntas de cierre}

\begin{enumerate}[label=\arabic*.]
    \item ¿Por qué medir mal la variable dependiente afecta la precisión?
    \item ¿Por qué medir mal la variable independiente genera sesgo de atenuación?
    \item ¿Qué significa que el coeficiente se acerque a cero?
    \item ¿Cómo influye la varianza del error de medición en la magnitud del sesgo?
\end{enumerate}

\end{document}
